\section*{Chapter \ref{ch:b2:aromatic-substitution} --- Aromaticity and Electrophilic Aromatic Substitution}

\begin{solution}{exo:b2:aromatic-substitution:1}
\ce{HNO3 + 2H2SO4 -> NO2+ + H3O+ + 2HSO4-}: nitric acid is protonated, then loses water.
\end{solution}

\begin{solution}{exo:b2:aromatic-substitution:2}
Toluene gives 2-bromotoluene and 4-bromotoluene: the methyl directs to the ortho and
para positions. Nitrobenzene gives 3-bromonitrobenzene, more slowly: the nitro group
directs meta.
\end{solution}

\begin{solution}{exo:b2:aromatic-substitution:3}
\ce{-CH3}: activating, o/p. \ce{-OH}: strongly activating, o/p. \ce{-Cl}: deactivating, o/p.
\ce{-NO2}: strongly deactivating, meta. \ce{-COCH3}: deactivating, meta. \ce{-NHCOCH3}:
activating (moderately), o/p.
\end{solution}

\begin{solution}{exo:b2:aromatic-substitution:4}
Acetophenone (1-phenylethan-1-one), \ce{C6H5COCH3}, through the \omterm{def:b2:aromatic-substitution:friedel-crafts}{acylium ion} \ce{CH3CO+}.
\end{solution}

\begin{solution}{exo:b2:aromatic-substitution:5}
The primary carbocation (or its complex) rearranges by a hydride shift to the secondary
isopropyl cation, which alkylates the ring. Propylbenzene: \omterm{def:b2:aromatic-substitution:friedel-crafts}{Friedel--Crafts acylation} with
propanoyl chloride (no rearrangement), then reduction of the ketone \ce{C=O} to \ce{CH2}.
\end{solution}

\begin{solution}{exo:b2:aromatic-substitution:6}
The \ce{-OH} group is a strong mesomeric donor: the ring is so activated that molecular bromine
is a sufficient electrophile, and all three ortho and para positions are substituted.
\end{solution}

\begin{solution}{exo:b2:aromatic-substitution:7}
3-Bromonitrobenzene: nitrate, then brominate (nitro directs meta). 4-Bromonitrobenzene:
brominate, then nitrate (bromine directs ortho/para), and separate the para isomer from the
ortho one.
\end{solution}

\begin{solution}{exo:b2:aromatic-substitution:8}
The nitrating mixture oxidises the very electron-rich aniline, and protonates its nitrogen:
\ce{-NH3+} is a meta-directing, \omterm{def:b2:aromatic-substitution:directing}{deactivating group}. In acetanilide the nitrogen lone pair is
shared with the carbonyl: less basic and less activating, it is not protonated or oxidised,
and still directs ortho/para, mainly para for steric reasons.
\end{solution}

\begin{solution}{exo:b2:aromatic-substitution:9}
Sulfonate toluene (mainly para, the bulky group avoids ortho); brominate: the bromine enters
ortho to the methyl (the para position is blocked, and both groups favour that position);
remove the \ce{SO3H} group by heating in dilute acid.
\end{solution}

\begin{solution}{exo:b2:aromatic-substitution:10}
Both groups direct ortho/para; their para positions are occupied by each other. Methoxy, the
stronger activator, wins: nitration ortho to \ce{-OCH3} (position 2), giving
4-methyl-2-nitroanisole.
\end{solution}

\begin{solution}{exo:b2:aromatic-substitution:11}
Each nitro group strongly deactivates the ring: the second nitration needs stronger conditions
than the first, the third (on a ring bearing two nitro groups and only the weakly activating
methyl) much stronger still.
\end{solution}

\begin{solution}{exo:b2:aromatic-substitution:12}
4-Nitrobenzoic acid: nitrate toluene, separate the para isomer, oxidise the methyl to
\ce{-COOH} (hot permanganate). 3-Nitrobenzoic acid: oxidise toluene to benzoic acid first, then
nitrate (\ce{-COOH} directs meta).
\end{solution}

\begin{solution}{pb:b2:aromatic-substitution:1}
\textbf{1.} \ce{HNO3 + 2H2SO4 -> NO2+ + H3O+ + 2HSO4-}.
\textbf{2.} \ce{O=N+=O}: linear, isoelectronic with \ce{CO2}.
\textbf{3.} The nitrogen bonds to the ring carbon para to \ce{CH3}, which becomes tetrahedral
(bearing H and \ce{NO2}).
\textbf{4.} The positive charge on the two carbons ortho to the attacked carbon and on the
carbon para to it, which bears the methyl.
\textbf{5.} The first, formation of the \omterm{def:b2:aromatic-substitution:seAr}{Wheland intermediate}.
\textbf{6.} A base of the medium: \ce{HSO4-} or water.
\textbf{7.} Two ortho, two meta, one para.
\textbf{8.} 2 : 2 : 1, that is 40~\% ortho, 40~\% meta, 20~\% para.
\textbf{9.} For meta attack, the positive charge is never on the carbon bearing the methyl; no
structure is helped by the methyl.
\textbf{10.} One resonance structure has the charge on the carbon bearing \ce{CH3}: a tertiary
carbocation, stabilised by the inductive effect and \omterm{def:b2:fragment-orbitals:hyperconjugation}{hyperconjugation} of the methyl.
\textbf{11.} Ortho and para are favoured (96~\% against 60~\% at random); meta is almost absent.
\textbf{12.} Per position: ortho $58/2 = 29~\%$, para 38~\%: para is the more reactive position,
the ortho positions being slightly hindered by the methyl.
\textbf{13.} Faster: the methyl activates the ring.
\textbf{14.} 4-Nitrotoluene (melts at \qty{52}{\degreeCelsius}).
\textbf{15.} Cool the crude mixture: the para isomer crystallises and is filtered off; the
liquid keeps the ortho and meta isomers.
\textbf{16.} A symmetric molecule packs better in a crystal: more interactions per volume,
a higher melting point.
\textbf{17.} By \omterm{def:b2:liquid-vapour-diagrams:fractional-distillation}{fractional distillation} under reduced pressure, or by chromatography; or by
further cooling (the meta isomer melts at \qty{16}{\degreeCelsius}).
\textbf{18.} The isomers have close boiling points (same formula, similar polarity).
\textbf{19.} \ce{C7H8}: \qty{92.14}{g/mol}; \ce{C7H7NO2}: \qty{137.14}{g/mol}.
\textbf{20.} $100/92.14 = \qty{1.085}{mol}$ of toluene; $0.90 \times 1.085 = \qty{0.977}{mol}$ of
nitrotoluenes.
\textbf{21.} Total $0.977 \times 137.14 = \qty{134.0}{g}$: ortho \qty{77.7}{g}, meta
\qty{5.4}{g}, para \qty{50.9}{g}.
\textbf{22.} A second nitration (dinitrotoluenes), and oxidation of the methyl group.
\textbf{23.} Oxidise the methyl to \ce{-COOH} with hot alkaline permanganate, then acidify.
\textbf{24.} The run gives $\boldsymbol{\approx \qty{51}{g}}$ of 4-nitrotoluene.
\end{solution}
